\section{Split designs}
\label{si:splits}

Every split keeps all oxygen sites of one host oxide on the same side, so no chemically identical compound is
seen in training and in testing; in this database one formula is one host, so grouping by formula and by host
coincide. The data that enter the splits are those of the filter cascade (Figure~\ref{fig:cascade}; counts in
Table~\ref{tab:cascade}), and the four designs are drawn in Figure~\ref{fig:splits}.

\begin{figure}[htbp]
\centering
\sifig{fig_cascade}
\caption{Filter cascade that defines universe v1. From the completed neutral oxygen-vacancy entries of the
release of Kumagai et al., the cascade keeps their neutral machine-learning subset, which already excludes the
perturbed-host-state entries, the entries with undetermined band edges and the dynamically unstable hosts. Three
filters follow in order: D1 removes the entries flagged as a split vacancy, an unknown defect type or a
configuration different from the initial one; D2 removes the host Sr$_2$Zr$_7$O$_{16}$, whose site labels do not
match the targets of the release; D3 keeps an entry of the two hosts with inconsistent cell metadata only if its
vacancy site maps to an oxygen atom with a matching first cation shell, and removes nothing. Dashed boxes are the
entries excluded at each step. Site and host counts after every step are in Table~\ref{tab:cascade}.}
\label{fig:cascade}
\end{figure}

In the primary protocol each of \numCfgSplitNOuterResamples{} outer resamples holds out
\numSplitMetaNTestHosts{} test hosts, a fraction \numSplitMetaTestFraction{} of the universe, stratified by the
host-mean formation energy, and orders the remaining pool of \numSplitMetaPoolSize{} hosts so that the training
set at budget $B$ is its first $B$ hosts. The budgets, \numCfgBudgetsHostsItemOne{} to \numCfgBudgetsHostsItemSix{}
hosts, are therefore nested within a resample, and the test hosts of a resample are the same at every budget, so
that the learning curves compare budgets on identical test oxides. Across the resamples,
\numSplitMetaCoverageInTestAtLeastOnce{} hosts are tested at least once and \numSplitMetaCoverageNeverInTest{}
never. The precision extension adds resamples drawn by the same rule, to \numPrecisionABFourHundredBSixHundredFiftyFourSubsetsThirtyXFiveNResamples{}
resamples with \numPrecisionABFourHundredBSixHundredFiftyFourSubsetsThirtyXFiveNSeeds{} training seeds each at
the two largest budgets (extension).

The published split of Kiyohara et al.\ \citep{Kiyohara2025machinelearning} is used once, to compare the
structure-only network with the published accuracy: the networks train on its
\numFilterCascadeKiyoharaSplitHostsTrain{} training hosts, stop early on its \numFilterCascadeKiyoharaSplitHostsVal{}
validation hosts and are tested on its \numFilterCascadeKiyoharaSplitHostsTest{} test hosts
(\numFilterCascadeKiyoharaSplitSitesTest{} sites). This comparison is not fully held out at the level of the
settings. The tuning used the nested training sets of resample \numCfgTuningAnchorResample{}, which was drawn
independently of the published split, so \numCZeroOfficialTuningOverlapWithKiyoharaTestSixHundredFiftyFourInKiyoharaTest{}
of the published test hosts lie in the largest tuning anchor
(\numCZeroOfficialTuningOverlapWithKiyoharaTestSixHundredFiftyFourOfWhichValidationHosts{} of them among its
validation hosts), and the comparison uses the settings of that anchor. The network weights never saw the labels
of those hosts in the comparison itself (\texttt{results/c0\_official.json}, key
\texttt{tuning\_overlap\_with\_kiyohara\_test}).

The grouped $K$-fold cross-validation (extension) partitions all \numSplitCvNHosts{} hosts into
$K = \numSplitCvKsItemOne{}$ or $K = \numSplitCvKsItemTwo{}$ folds, so that every host is tested exactly once and the
mean training set, \numCvGOneVOneXTrainingHostsKTen{} and \numCvGOneVOneXTrainingHostsKTwenty{} hosts, exceeds the
largest registered budget. Its test population is the whole universe, including the
\numSplitMetaCoverageNeverInTest{} hosts that no outer resample tests, so a cross-validated error is not directly
comparable with a registered budget at a similar number of training hosts.

Leave-one-chemistry-out asks whether the models transfer to chemistries absent from training. Each host belongs to
the family given by the set of periodic-table groups of its cations, which gives \numSplitLocoNFamilies{} families
over the universe, and \numSplitLocoNFolds{} folds each hold out whole families. The runs of the originally
registered network follow the registered design but are the collaborator's, whose validation carve-out differs
from the one our later clarification specifies; they are adopted as run and labelled extension at the results
freeze (\texttt{docs/freeze\_notes.md}), while \texttt{docs/provenance\_labels.md} records them as registered in
design and leaves the choice to the author. The runs of the primary network are post hoc.

\begin{figure}[htbp]
\centering
\sifig{fig_splits}
\caption{Split designs, drawn as sets of hosts (dark: test; light: training; mid-grey: validation). Every split
keeps all sites of a host on one side. Widths are schematic, not to scale. (\textbf{a}) Primary protocol: each
outer resample holds out \numSplitMetaNTestHosts{} test hosts, stratified by host-mean $E_\mathrm{f}$, and orders
the remaining pool so that the training set at budget $B$ is the first $B$ hosts; budgets are nested, and the
learning curves are drawn against $B$. Validation hosts are carved out of each training set. (\textbf{b}) The
published split of Kiyohara et al.: the networks train on its training hosts, stop early on its validation hosts
and are tested on its test hosts. (\textbf{c}) Grouped $K$-fold cross-validation ($K = \numSplitCvKsItemOne{}$ and
$\numSplitCvKsItemTwo{}$; extension): every host is tested exactly once, with training sets larger than the largest
budget of (a). (\textbf{d}) Leave-one-chemistry-out: hosts are grouped into chemistry families (the set of
periodic-table groups of the cations), and \numSplitLocoNFolds{} folds, each holding out whole families, test
transfer to chemistries absent from training.}
\label{fig:splits}
\end{figure}
\FloatBarrier
